\clearpage
\VerdierPhysicalPageAnchor{0197}%
\begingroup\small\setlength{\parskip}{0pt}\setlength{\emergencystretch}{1.5em}
\setlength{\abovedisplayskip}{4pt plus 2pt minus 2pt}%
\setlength{\belowdisplayskip}{4pt plus 2pt minus 2pt}%

\VerdierSemanticAnchor{ch3d-2-3-1-induction-step-conclusion}
そこで $L_q^\bullet=L_{q-1}^\bullet(\alpha,-q)$、
$w_q=\varphi$、$v_q=p(\alpha,-q)$ とおく。複体 $L_n^\bullet$ と複体の射
$w_n$、$v_n$ は、$n\leq q$ に対して性質 2) と 3) をもつ。あとは
$L_q^\bullet$ が上の性質 1) をもつことを示せばよい。これは集合
$M^\wedge$ が完全、すなわち $(E_6)$（2.2.4）を満たすことから従う。

\VerdierSemanticAnchor{ch3d-2-3-1-proof-b}
b) を証明する。$L$ を $M$ に属する対象とし、
$\dot u\colon X\longrightarrow L$ をエピ射とする。$Y$ を
$\dot u$ の核とおく。次の擬同型がある：

\VerdierSemanticAnchor{ch3d-2-3-1-kernel-quasi-isomorphism}
\begin{center}
\begin{tikzcd}[column sep=2.25em,row sep=2.6em]
\cdots \arrow[r] & 0 \arrow[r] \arrow[d] & Y \arrow[r] \arrow[d]
  & X \arrow[r] \arrow[d,"\dot u"] & 0 \arrow[r] \arrow[d] & \cdots \\
\cdots \arrow[r] & 0 \arrow[r] & 0 \arrow[r] & L \arrow[r] & 0 \arrow[r]
  & \cdots.
\end{tikzcd}
\end{center}
したがって、次の擬同型がある：

\VerdierSemanticAnchor{ch3d-2-3-1-resolution-quasi-isomorphism}
\begin{center}
\begin{tikzcd}[column sep=2.25em,row sep=2.6em]
\cdots \arrow[r] & L^{-1} \arrow[r] \arrow[d] & L^0 \arrow[r] \arrow[d,"v"]
  & L^1 \arrow[r] \arrow[d] & L^2 \arrow[r] \arrow[d] & \cdots \\
\cdots \arrow[r] & Y \arrow[r] & X \arrow[r] & 0 \arrow[r] & 0 \arrow[r]
  & \cdots,
\end{tikzcd}
\end{center}

\VerdierSemanticAnchor{ch3d-2-3-1-e2}
ここで各 $L^i$ は $M$ に属する。したがって射 $\dot u v$ はエピ射であり、
性質 $(E_2)$ が示された。次に、

\VerdierSemanticAnchor{ch3d-2-3-1-exact-resolution}
\[
0\longrightarrow X\xrightarrow{\dot u}L_n\longrightarrow L_{n-1}
\longrightarrow\cdots\longrightarrow L_0\longrightarrow0
\]
を完全列とし、各 $L_i$ は $M$ に属するとする。このとき次の擬同型がある：
\begin{center}
\begin{tikzcd}[column sep=1.35em,row sep=2.6em]
\cdots \arrow[r] & 0 \arrow[r] \arrow[d] & 0 \arrow[r] \arrow[d]
  & X \arrow[r] \arrow[d,"\dot u"] & 0 \arrow[r] \arrow[d]
  & \cdots \arrow[r] & 0 \arrow[r] \arrow[d] & 0 \arrow[r] \arrow[d]
  & \cdots \\
\cdots \arrow[r] & 0 \arrow[r] & 0 \arrow[r] & L_n \arrow[r]
  & L_{n-1} \arrow[r] & \cdots \arrow[r] & L_0 \arrow[r] & 0 \arrow[r]
  & \cdots.
\end{tikzcd}
\end{center}
従って、次の擬同型が存在する：

\VerdierSemanticAnchor{ch3d-2-3-1-e3-quasi-isomorphism}
\begin{center}
\begin{tikzcd}[column sep=2.7em,row sep=2.6em]
\cdots \arrow[r] & L^{-n-1} \arrow[r] \arrow[d]
  & L^{-n} \arrow[r] \arrow[d,"v"] & L^{-n+1} \arrow[r] \arrow[d]
  & \cdots \\
\cdots \arrow[r] & 0 \arrow[r] & X \arrow[r] & 0 \arrow[r] & \cdots
\end{tikzcd}
\end{center}
\endgroup
